% Chapter 1: A Century of Experiments
% Subsections in chronological order, each a key experiment or result

To aid comparison across these different traditions, this chapter adopts a single, modern set of notation for the core concepts of reinforcement learning, following current conventions \citep{sutton2018}. We denote a state by $s$, an action by $a$, a policy by $\pi$, the state-value function by $V(s)$, the action-value function by $Q(s,a)$, the reward by $r$, and the discount factor by $\gamma$. When presenting historical formalisms, we will often use the author's original notation first, and then explicitly map it to this modern framework. This sometimes requires a slight abuse of notation, as the original concepts may not be perfect analogues, but it serves the goal of creating a unified conceptual narrative.

\subsection{Thorndike's Cats (1898)}

Comparative psychology in the 1890s sought to measure animal intelligence through controlled experiment rather than anecdote. \citet{Thorndike1898} placed hungry cats inside latched wooden boxes with food visible outside and recorded escape latency across repeated trials.

Trial 1 required 160 seconds of undirected clawing, biting, and squeezing through slats. By trial 24 the cat pressed the latch directly in 7 seconds. The learning curve was smooth rather than stepwise, consistent with incremental strengthening of stimulus-response associations rather than sudden insight.\footnote{Thorndike's escape-time data followed an approximate power law $T(n) \propto n^{-\beta}$ rather than an exponential decay, a pattern later termed the ``power law of practice.'' Whether animal and human learning curves are fundamentally power-law or a mixture of exponentials remains debated in the quantitative learning literature.} Thorndike codified this pattern as the Law of Effect:

\begin{quote}\itshape
Of several responses made to the same situation, those which are accompanied or closely followed by satisfaction to the animal will, other things being equal, be more firmly connected with the situation, so that, when it recurs, they will be more likely to recur \citep{Thorndike1898}.
\end{quote}

In modern terms, the box configuration is the state $s$, the motor action (pressing, pulling, clawing) is $a$, and escape followed by food is the reward $r$.\footnote{This is a slight abuse of notation, as Thorndike's ``situation'' was not a formally defined state and his ``responses'' were not drawn from a well-defined action set. The mapping is conceptual, highlighting the analogy between the Law of Effect and the principles of modern reinforcement learning.} The ``stamping in'' of successful responses increases the probability of selecting that action given the state, a directional update rule that Q-learning would formalize nine decades later.

\subsection{Samuel's Checkers Program (1959)}

\citet{Samuel1959} set out to build a checkers program for the IBM 704 that could improve through experience rather than manual tuning, reasoning that learning from play should eventually reduce the burden of hand-coding evaluation logic. Samuel's deeper bet was that learning from experience could eventually surpass hand-coded knowledge entirely, a claim that TD-Gammon would vindicate three decades later and that AlphaGo Zero would prove definitively by exceeding all human and computer Go programs without any human training data. The program learned through self-play using two complementary mechanisms.

Rote learning stored every board position encountered along with its backed-up value, creating an expanding lookup table over the course of play. Generalization learning evaluated positions with a linear function
\begin{equation}
V(s) = \sum_{i} w_i \phi_i(s)
\end{equation}
where $\phi_i(s)$ were hand-crafted features: piece advantage, piece advancement, center control, mobility, and king-row threat. After each move, the weights were updated according to
\begin{equation}
\mathbf{w} \leftarrow \mathbf{w} + \alpha \left[ V(s_{t+1}) - V(s_t) \right] \nabla_{\mathbf{w}} V(s_t)
\end{equation}
The program reached expert-level play by 1962 and defeated Robert Nealey, a self-proclaimed expert, in a publicized match. Samuel's update is precisely TD(0) with linear function approximation, a connection formalized by \citet{sutton1988} twenty-nine years later.\footnote{Alpha-beta pruning, implemented following McCarthy's suggestion, reduces the minimax search complexity from $O(b^d)$ to $O(b^{d/2})$ under optimal move ordering, where $b$ is the branching factor and $d$ is the search depth. For checkers ($b \approx 8$, $d = 10$), this reduces the search space from $8^{10} \approx 10^9$ to $8^5 \approx 33{,}000$ positions.}

The linear polynomial assumed that features contributed independently to position value. In a follow-up study, Samuel replaced it with \emph{signature tables}: hierarchical lookup tables that captured within-group feature interactions. The 27 board parameters were divided into 9 groups of 3, and the evaluation became
\begin{equation}
V(s) = F\!\Bigl(T_1\bigl(\phi_1(s), \phi_2(s), \phi_3(s)\bigr),\; \ldots,\; T_9\bigl(\phi_{25}(s), \phi_{26}(s), \phi_{27}(s)\bigr)\Bigr)
\end{equation}
where each $T_g$ is a lookup table mapping a discretized feature triple to a score, and $F$ combines the group outputs through two further levels of table lookup. With parameters restricted to 3, 5, or 7 discrete values, the full three-level hierarchy required 20{,}956 table entries.\footnote{Each group of 3 parameters, discretized to $k \in \{3, 5, 7\}$ levels, yields a table of $k^3$ entries. With two further levels of table lookup combining group outputs, the total capacity was 20{,}956 entries stored in the IBM 7094's memory. This hierarchical decomposition can be viewed as a factored function approximation: it captures within-group interactions (up to 3-way) while assuming independence across groups.} This replacement of a linear combination with a hierarchical nonlinear mapping is a direct precursor to neural network function approximation: the conceptual step from Samuel's tables to Tesauro's hidden layers is replacing discrete lookups with continuous, differentiable activations.

Both the polynomial coefficients and the table entries were learned from a library of master-level games.\footnote{The library contained approximately 250{,}000 positions extracted from published master games. Each position was paired with the book-recommended move, providing a supervised training signal. By modern standards this is a small dataset; contemporary chess engines train on billions of self-play positions.} For each board parameter $i$, the program counted the number $H_i$ of non-book moves for which the parameter value exceeded that of the book-recommended move, and the number $L_i$ for which it fell below. The learned coefficient was
\begin{equation}
C_i = \frac{H_i - L_i}{H_i + L_i}
\end{equation}
This book learning procedure is supervised learning from expert demonstrations, not the self-play bootstrapping described above. Samuel used both: book learning to initialize coefficients, self-play to refine them. The distinction between learning from expert data and learning from self-generated experience anticipates the modern contrast between imitation learning and reinforcement learning.

\citet{Minsky1961} organized artificial intelligence into five problem areas (search, pattern recognition, learning, planning, and induction) in his contemporary survey. Among the obstacles he identified was the ``mesa phenomenon'': local optima in hill-climbing search that appear optimal from nearby but are not globally optimal, a structure that foreshadows the exploration-exploitation dilemma. He also identified the credit assignment problem: in a game lasting many moves, which of the early decisions produced the eventual win or loss? Samuel's temporal difference update is one answer to this question, propagating information about outcomes backward through the sequence of positions without waiting for the game to end.

\subsection{The Rescorla-Wagner Model (1972)}

Classical conditioning theory through the mid-twentieth century held that temporal pairing of stimuli was sufficient for associative learning. The blocking experiment provided decisive evidence against this view, showing that prediction error, not mere co-occurrence, governs associative change.

\begin{quote}\itshape
Organisms only learn when events violate their expectations. Certain expectations are built up about the events following a stimulus complex; expectations initiated by that complex and its component stimuli are then only modified when consequent events disagree with the composite expectation \citep{RescorlaWagner1972}.
\end{quote}

The protocol proceeds in two phases. In Phase I, a rat learns that a tone predicts shock and develops a conditioned fear response to the tone alone. In Phase II, the compound stimulus of tone plus light is paired with the same shock. When the light is subsequently presented alone, no fear response occurs. The light is ``blocked'' because the tone already predicts the shock perfectly, leaving no prediction error to drive learning about the light. \citet{RescorlaWagner1972} formalized this observation with the following update rule:
\begin{equation}
\Delta V_i = \alpha_i \beta_j (\lambda_j - V_{\text{tot}})
\end{equation}
where $V_i$ is the associative strength of stimulus $i$, $\alpha_i$ is the salience of stimulus $i$, $\beta_j$ is the learning rate associated with unconditioned stimulus $j$, $\lambda_j$ is the maximum conditioning that stimulus $j$ supports, and $V_{\text{tot}} = \sum_k V_k$ is the total prediction from all stimuli present on the trial.\footnote{The Rescorla-Wagner update $\Delta V_i = \alpha_i \beta_j (\lambda_j - V_{\text{tot}})$ is mathematically identical to the Widrow-Hoff (1960) least mean squares (LMS) rule used in adaptive signal processing. Both minimize mean squared prediction error by stochastic gradient descent. The parallel development in psychology and engineering was independent.} Learning halts when $V_{\text{tot}} = \lambda_j$, because the prediction error $(\lambda_j - V_{\text{tot}})$ reaches zero. In the blocking design, $V_{\text{tone}}$ already equals $\lambda$ after Phase I, so $\Delta V_{\text{light}} = 0$ throughout Phase II.

The Rescorla-Wagner model can be mapped directly to the temporal difference framework. The term $V_{\text{tot}}$ is the value of the current state, $V(s)$. The term $\lambda_j$ represents the maximum possible reward, which in a single-step episode is the reward itself, $r$. The associative strength of a single stimulus, $V_i$, corresponds to a component of the value function. With a discount factor $\gamma=0$, the update is driven by the prediction error $(\lambda_j - V_{\text{tot}})$, which is equivalent to the TD error, $\delta = r - V(s)$. The model is thus a special case of TD learning for single-step, terminal-reward tasks.

\subsection{Temporal Difference Learning and Q-Learning (1988, 1989)}

The gap between dynamic programming, which requires a complete model of the environment, and Monte Carlo estimation, which requires waiting for complete episodes, motivated a method that could learn incrementally from partial experience.

\begin{quote}\itshape
Whereas conventional prediction-learning methods are driven by the error between predicted and actual outcomes, TD methods are similarly driven by the error or difference between temporally successive predictions; with them, learning occurs whenever there is a change in prediction over time \citep{sutton1988}.
\end{quote}

\citet{sutton1988} formalized temporal difference learning and demonstrated its properties on a five-state random walk.

States A through E are arranged in a line with absorbing terminal states at each end. At each step the agent moves left or right with equal probability. The right terminal yields reward 1 and the left terminal yields reward 0. The true state values are $\frac{1}{6}, \frac{2}{6}, \frac{3}{6}, \frac{4}{6}, \frac{5}{6}$. Sutton compared TD(0), which updates after every transition, with Monte Carlo estimation, which updates only at the end of each episode. TD(0) converged faster and with less data than Monte Carlo on this problem.

The general TD($\lambda$) update combines both approaches through an eligibility trace:
\begin{equation}
V(s_t) \leftarrow V(s_t) + \alpha \, \delta_t \, e_t(s)
\end{equation}
where $\delta_t = r_{t+1} + \gamma V(s_{t+1}) - V(s_t)$ is the TD error and $e_t(s) = \gamma \lambda \, e_{t-1}(s) + \mathbbm{1}\{s = s_t\}$ is the eligibility trace for state $s$.\footnote{The eligibility trace formulation (backward view) is computationally efficient, requiring only $O(|\mathcal{S}|)$ storage per step. The equivalent forward view computes the $\lambda$-return $G_t^\lambda = (1-\lambda) \sum_{n=1}^{\infty} \lambda^{n-1} G_t^{(n)}$, a geometrically weighted average of $n$-step returns. The backward and forward views produce identical updates in the offline (batch) case; they differ slightly online due to the accumulating trace approximation.} Intermediate values of $\lambda$ interpolate between pure bootstrapping ($\lambda = 0$) and Monte Carlo returns ($\lambda = 1$), and the optimal $\lambda$ is problem-dependent. Updating from the next estimate rather than waiting for the final outcome allows learning to proceed at every timestep, even in continuing tasks where Monte Carlo methods cannot be used.

TD methods as developed by Sutton solve the prediction problem: estimating the value of a given policy. The control problem, finding the optimal policy without access to the transition dynamics, requires a different approach. \citet{WatkinsDayan1992}, formalizing Watkins's 1989 PhD thesis, provided an incremental method for dynamic programming that imposes limited computational demands:
\begin{equation}
Q(s_t, a_t) \leftarrow Q(s_t, a_t) + \alpha \left[ r_{t+1} + \gamma \max_{a'} Q(s_{t+1}, a') - Q(s_t, a_t) \right]
\end{equation}
The maximization over $a'$ makes Q-learning off-policy: the update target uses the greedy action at the next state regardless of the action actually taken, so the agent can explore with an $\varepsilon$-greedy behavior policy while learning the optimal policy directly.\footnote{The $\max$ operator in the Q-learning target introduces a positive bias: $\mathbb{E}[\max_a Q(s,a)] \geq \max_a \mathbb{E}[Q(s,a)]$ when Q-values are noisy estimates. This maximization bias can cause systematic overestimation of action values. Double Q-learning (van Hasselt, 2010) addresses this by maintaining two independent Q-functions and using one to select and the other to evaluate, decoupling the max from the estimate.}

Convergence to $Q^*$ in the tabular setting requires two conditions: every state-action pair must be visited infinitely often, and the learning rate must satisfy the standard stochastic approximation conditions.\footnote{The conditions $\sum_t \alpha_t = \infty$ and $\sum_t \alpha_t^2 < \infty$ are the Robbins-Monro (1951) stochastic approximation conditions. The first ensures the step sizes are large enough in aggregate to reach the optimum from any initialization. The second ensures they shrink fast enough for the noise to average out. A common schedule satisfying both is $\alpha_t = c / (c + t)$ for constant $c > 0$. In practice, fixed small learning rates are often used instead, which do not guarantee convergence but track non-stationarity.} Under these conditions, $Q(s,a) \to Q^*(s,a)$ for all $(s,a)$ with probability one. The optimal policy is then $\pi^*(s) = \argmax_a Q^*(s,a)$.

The off-policy advantage is cleanly demonstrated by the cliff-walking problem \citep{sutton2018}. A $4 \times 12$ gridworld has a start cell in the bottom-left corner and a goal cell in the bottom-right corner, with a cliff running along the entire bottom row between them. Stepping onto the cliff yields a reward of $-100$ and returns the agent to start; all other steps yield $-1$. Q-learning, updating toward the greedy $\max_{a'} Q(s_{t+1}, a')$, learns the optimal shortest path directly along the cliff edge, because its target policy is greedy regardless of the exploratory actions actually taken. SARSA, which updates toward the action actually selected under the $\varepsilon$-greedy behavior policy, learns a longer safe path one row above the cliff, because its on-policy updates incorporate the cost of occasional exploratory steps off the edge. The two algorithms converge to different policies under the same $\varepsilon$-greedy exploration, isolating the consequence of off-policy versus on-policy learning.

\subsection{Williams' REINFORCE (1992)}

Value-based methods learn an action-value function and derive a policy from it. An alternative approach optimizes the policy directly by gradient ascent on expected return. \citet{williams1992} derived the policy gradient for a parameterized stochastic policy $\pi_\theta(a|s)$:
\begin{equation}
\nabla_\theta J(\theta) = \mathbb{E}_{\pi_\theta} \left[ \sum_{t=0}^{T} \nabla_\theta \log \pi_\theta(a_t|s_t) \, G_t \right]
\end{equation}
where $G_t = \sum_{k=0}^{T-t} \gamma^k r_{t+k+1}$ is the discounted return from time $t$. The log-derivative trick, $\nabla_\theta \pi_\theta = \pi_\theta \nabla_\theta \log \pi_\theta$, allows the gradient to be estimated from sampled trajectories without differentiating through the environment dynamics or the transition model.

The practical difficulty is variance. Because $G_t$ is a sum of stochastic rewards over the remainder of the episode, and because different trajectories may yield very different returns, the gradient estimates have high variance. Variance reduction techniques, including baselines $b(s_t)$ subtracted from $G_t$ without introducing bias, are essential for practical performance.\footnote{The variance-minimizing baseline for the REINFORCE estimator is $b^*(s_t) = \mathbb{E}[G_t^2 \|\nabla_\theta \log \pi_\theta\|^2] / \mathbb{E}[\|\nabla_\theta \log \pi_\theta\|^2]$, which depends on the gradient magnitudes. In practice, $b(s_t) \approx V^\pi(s_t)$ is used, leading to the advantage actor-critic architecture where $G_t - b(s_t) \approx A^\pi(s_t, a_t)$.}

To see what this buys in practice, consider a robot arm that must apply a continuous torque $a \in \mathbb{R}$ to reach a target angle. Q-learning requires computing $\max_a Q(s,a)$ at every update, which is tractable when $a$ comes from a finite set but becomes a nested optimization problem when the action space is continuous. A policy gradient method sidesteps the issue entirely: parameterize the policy as a Gaussian $\pi_\theta(a|s) = \mathcal{N}(\mu_\theta(s),\, \sigma_\theta^2(s))$, sample an action, observe the return, and update $\theta$ by REINFORCE. The gradient $\nabla_\theta \log \pi_\theta(a|s)$ is available in closed form for the Gaussian, and no maximization over actions is required. This is one reason that most continuous-control results in reinforcement learning, from simulated locomotion to real robotic manipulation, use the policy gradient framework rather than Q-learning.

Policy gradients also handle problems where the optimal policy is genuinely stochastic. In a partially observed environment where distinct underlying states produce the same observation, a deterministic policy must take the same action in both states, which may be suboptimal. A stochastic policy can mix over actions and achieve higher expected return. \citet{SuttonMcAllester2000} emphasized this advantage, noting that the value-function approach is ``oriented toward finding deterministic policies, whereas the optimal policy is often stochastic.'' They also identified a subtler problem with value-based methods: an arbitrarily small change in estimated Q-values can cause the greedy policy to switch actions discontinuously, destabilizing learning with function approximation. Policy gradients avoid this because small changes in $\theta$ produce small changes in $\pi_\theta$.

\citet{SuttonMcAllester2000} demonstrated this concretely with a short corridor of four states where, in one interior state, the effect of the two actions is reversed: ``left'' moves right and ``right'' moves left. Any deterministic policy either always moves in the correct direction or always moves in the wrong direction at that state, and neither achieves the optimal expected return. A stochastic policy that mixes left and right with the correct probabilities does strictly better. An $\varepsilon$-greedy value-based method cannot discover this mixture because its policy is deterministic except for uniform random noise, whereas a policy gradient method parameterizes the action probabilities directly and converges to the optimal stochastic policy by gradient ascent on expected return.

\citet{SuttonMcAllester2000} proved the policy gradient theorem in the average reward setting, establishing that the gradient of expected return with respect to policy parameters takes the form $\nabla_\theta J(\theta) = \mathbb{E}_{s \sim d^\pi, a \sim \pi_\theta} [Q^\pi(s,a) \nabla_\theta \log \pi_\theta(a|s)]$, where $d^\pi$ is the stationary state distribution under $\pi$. The result is remarkable for what it excludes: the gradient does not require differentiating through $d^\pi$, even though $d^\pi$ depends on $\theta$ through the policy's effect on state transitions. The policy gradient theorem enabled the development of actor-critic methods that combine the low variance of value function estimation with the generality of policy parameterization.

\subsection{TD-Gammon (1994)}

The first large-scale demonstration that temporal difference learning with neural network function approximation could produce expert-level play was \citet{tesauro1994}'s TD-Gammon. Backgammon has approximately $10^{20}$ legal positions, far beyond the reach of tabular methods or exhaustive search.

TD-Gammon trained a feedforward neural network to estimate the probability of winning from any board position. The input $\mathbf{x}(s) \in \mathbb{R}^{198}$ encoded the board position using a unary representation.\footnote{The encoding used four binary units per board point per player: for each of the 24 points and each player, 4 binary units indicated occupancy of 1, 2, 3, or $\geq 4$ pieces. This yields $24 \times 2 \times 4 = 192$ units, plus 6 additional units for pieces on the bar and borne off, totaling 198. The unary encoding allows the network to learn nonlinear functions of piece count without requiring the hidden layer to first disentangle a scalar input.} A hidden layer of 80 sigmoid units (40 in version 2.0) fed a single sigmoid output:
\begin{equation}
\hat{V}(s) = \sigma\!\bigl(\mathbf{w}^\top \sigma(W\mathbf{x}(s) + \mathbf{b}) + c\bigr)
\end{equation}
where $\sigma$ is the logistic sigmoid, $W$ is the input-to-hidden weight matrix, and $\mathbf{w}$ is the hidden-to-output weight vector.

The network was trained by self-play using TD($\lambda$) with $\lambda = 0.7$. After each move from position $s_t$ to $s_{t+1}$, the weights $\boldsymbol{\theta}$ (comprising $W$, $\mathbf{w}$, $\mathbf{b}$, and $c$) were updated by
\begin{equation}
\boldsymbol{\theta} \leftarrow \boldsymbol{\theta} + \alpha \bigl[\hat{V}(s_{t+1}) - \hat{V}(s_t)\bigr] \, \mathbf{e}_t
\end{equation}
where the eligibility trace $\mathbf{e}_t = \sum_{k=1}^{t} \lambda^{t-k} \nabla_{\boldsymbol{\theta}} \hat{V}(s_k)$ accumulates exponentially decayed gradients of past predictions, computed by backpropagation. At game's end, $\hat{V}(s_{t+1})$ is replaced by the outcome $z \in \{0, 1\}$, grounding the bootstrapping chain in actual wins and losses. At each turn, the program selected the move maximizing $\hat{V}(\mathrm{next}(s_t, a))$ over legal actions, a purely greedy policy. The stochastic dice rolls in backgammon provided natural exploration, a property Tesauro conjectured was critical to the method's success.\footnote{The stochastic dice in backgammon provide forced exploration: different rolls lead to different positions regardless of the player's strategy. This means a greedy policy still visits a diverse set of states. Deterministic games like chess and Go lack this property, which partly explains why TD-Gammon's pure self-play approach did not immediately transfer to those domains. AlphaGo Zero later solved the exploration problem in Go through MCTS-based action selection with added Dirichlet noise at the root node.}

TD-Gammon 1.0, trained on 300{,}000 self-play games with 1-ply lookahead, played at a strong intermediate level. Version 2.1, trained on 1.5 million games with 2-ply lookahead and 80 hidden units, achieved near-parity with Bill Robertie, a former world champion, losing by a single point over a 40-game session. The program discovered novel positional strategies that were subsequently adopted by the human backgammon community.

\begin{quote}\itshape
From an a priori point of view, this methodology appeared unlikely to produce any sensible learning, because random strategy is exceedingly bad, and because the games end up taking an incredibly long time \ldots\ The rather surprising result, after tens of thousands of training games, was that a significant amount of learning actually took place \citep{tesauro1994}.
\end{quote}

\subsection{Baird's Counterexample and the Deadly Triad (1995)}

\citet{Baird1995} constructed a seven-state star MDP that demonstrated divergence of Q-learning with linear function approximation. The MDP has six ``outer'' states that all transition to a single ``inner'' state under the target policy. The off-policy behavior samples states uniformly. With linear function approximation, the weights grow without bound under repeated Q-learning updates, despite the simplicity of the MDP.

\begin{quote}\itshape
It is unfortunate that a reinforcement learning algorithm can be guaranteed to converge for lookup tables, yet be unstable for function-approximation systems that have even a small amount of generalization \citep{Baird1995}.
\end{quote}

The source of the instability is the interaction of three components: bootstrapping, off-policy learning, and function approximation. \citet{sutton2018} later named this combination the deadly triad. Convergence guarantees exist for any two of the three components in isolation; it is the simultaneous presence of all three that permits divergence.

This result explains the asymmetry between Tesauro's success and Baird's failure. TD-Gammon used bootstrapping and function approximation but was on-policy: the training data came from the same self-play policy whose value was being estimated. Baird's counterexample used all three components, including off-policy updates, and diverged.\footnote{The divergence in Baird's example arises because the TD fixed point minimizes the projected Bellman error $\|V_\theta - \Pi T V_\theta\|$, where $\Pi$ projects onto the function approximation space. When the projection and the Bellman operator $T$ interact adversarially under off-policy sampling, the fixed point may not exist or may be unstable. Gradient TD methods (Sutton et al., 2009) resolve this by performing stochastic gradient descent on the mean-squared projected Bellman error, guaranteeing convergence at the cost of a second set of learned weights.}

Baird did not only diagnose the problem but proposed a constructive solution: residual gradient algorithms that perform gradient descent on the mean-squared Bellman residual, guaranteeing convergence with function approximation at the cost of converging to a different fixed point than standard TD. This constructive direction connects to the Gradient TD methods mentioned above and to the stabilization techniques that DQN would later adopt in a different form: where Baird changed the learning rule to follow a true gradient, DQN changed the data distribution (replay) and the target (frozen network) to make the existing rule behave.

\subsection{Deep Q-Networks and TRPO/PPO (2015)}

The two decades after TD-Gammon saw no comparable success in applying neural networks to reinforcement learning. Attempts to replicate Tesauro's approach in other domains, including chess and Go, failed to produce competitive play, and the field largely retreated to linear methods or tabular representations. \citet{mnih2015} themselves noted this gap. End-to-end learning from raw sensory input became feasible when sufficient computation met architectural innovations for stabilizing the deadly triad. \citet{mnih2015} trained a single convolutional neural network to play 49 Atari 2600 games directly from pixel inputs ($210 \times 160 \times 3$) and a scalar score, using no game-specific features or hand-crafted representations. The network was trained by minimizing the squared temporal difference error:
\begin{equation}
L(\theta) = \mathbb{E}_{(s,a,r,s') \sim \mathcal{D}} \left[ \left( r + \gamma \max_{a'} Q(s', a'; \theta^-) - Q(s, a; \theta) \right)^2 \right]
\end{equation}
Two innovations stabilized learning. First, experience replay: transitions $(s, a, r, s')$ were stored in a replay buffer $\mathcal{D}$ and sampled uniformly for training, breaking the temporal correlation between consecutive updates. Second, a target network: the bootstrap target used a frozen copy of the parameters $\theta^-$, updated to match $\theta$ only every $C$ steps, so that the regression target does not shift with each gradient step.\footnote{The convolutional architecture used three layers (32 filters of $8 \times 8$ with stride 4, then 64 filters of $4 \times 4$ with stride 2, then 64 filters of $3 \times 3$ with stride 1) followed by a 512-unit fully connected layer. Input frames were downsampled to $84 \times 84$ grayscale and stacked 4 frames deep to provide velocity information. The replay buffer held $10^6$ transitions, the target network was updated every $C = 10{,}000$ steps, and training used RMSProp with a learning rate of $2.5 \times 10^{-4}$.}

The same architecture and hyperparameters exceeded human-level performance on 29 of 49 games, from Breakout to Space Invaders, establishing that a single network can learn diverse tasks from raw pixels without domain knowledge. Games requiring long-horizon planning or sparse rewards (such as Montezuma's Revenge) remained difficult, exposing the limits of value-based methods without hierarchical structure.

Policy gradient methods as described in the REINFORCE subsection suffer from a practical instability: a single large gradient step can move the policy into a region of parameter space where performance collapses, and recovery may be slow or impossible. The difficulty is that the objective $J(\theta)$ is not well approximated by its linear Taylor expansion when the step size is large relative to the curvature of the policy space. Gradient-free methods such as the cross-entropy method and random search were competitive with or outperforming gradient-based policy optimization on some continuous control tasks, suggesting that the optimization landscape, not the gradient signal itself, was the bottleneck. \citet{Schulman2015} addressed this problem with Trust Region Policy Optimization (TRPO), grounded in a monotonic improvement guarantee. For any two policies $\pi$ and $\tilde{\pi}$, the performance difference can be written exactly as
\begin{equation}
\eta(\tilde{\pi}) = \eta(\pi) + \mathbb{E}_{s \sim d^{\tilde{\pi}}, a \sim \tilde{\pi}} \left[ A^\pi(s,a) \right]
\end{equation}
where $\eta(\pi) = \mathbb{E}[\sum_t \gamma^t r_t]$ is the expected discounted return and $A^\pi(s,a) = Q^\pi(s,a) - V^\pi(s)$ is the advantage function. Because $d^{\tilde{\pi}}$ depends on the new policy and is difficult to compute, TRPO approximates it by using the surrogate objective
\begin{equation}
L_\pi(\tilde{\pi}) = \eta(\pi) + \mathbb{E}_{s \sim d^\pi, a \sim \tilde{\pi}} \left[ A^\pi(s,a) \right]
\end{equation}
which replaces $d^{\tilde{\pi}}$ with $d^\pi$. Schulman showed that the approximation error can be bounded:
\begin{equation}
\eta(\tilde{\pi}) \geq L_\pi(\tilde{\pi}) - C \, D_{\mathrm{KL}}^{\max}(\pi, \tilde{\pi})
\end{equation}
where $C$ is a constant depending on $\gamma$ and the advantage range,\footnote{The constant is $C = 4\gamma\epsilon / (1-\gamma)^2$ where $\epsilon = \max_{s,a} |A^\pi(s,a)|$, originally derived by Kakade and Langford (2002) for conservative policy iteration. The bound is loose: for $\gamma = 0.99$ and $\epsilon = 1$, $C \approx 400$, making the guaranteed improvement region extremely small. TRPO therefore ignores the bound and uses a fixed KL threshold $\delta \approx 0.01$, relying on empirical stability rather than the theoretical guarantee.} and $D_{\mathrm{KL}}^{\max}$ is the maximum KL divergence over states. Maximizing the right-hand side guarantees monotonic improvement. In practice, TRPO solves the constrained optimization problem
\begin{equation}
\max_\theta \; L_{\theta_{\text{old}}}(\theta) \quad \text{subject to} \quad \bar{D}_{\mathrm{KL}}(\theta_{\text{old}}, \theta) \leq \delta
\end{equation}
using the conjugate gradient method to compute the search direction and a line search to enforce the constraint. TRPO achieved strong results on continuous control tasks in MuJoCo (swimmer, hopper, walker) and on a subset of Atari games.

On the MuJoCo Hopper, the natural gradient method (fixed KL penalty) learned balanced standing but never walked forward, while TRPO learned a hopping gait. The failure mode is instructive: an oversized gradient step can shift the policy into a falling regime where collected data becomes uninformative and recovery is difficult; the KL constraint prevents this by bounding the policy change per update.

\citet{Schulman2017} proposed Proximal Policy Optimization (PPO) as a simpler alternative that achieves comparable performance without the conjugate gradient computation. PPO replaces the KL constraint with a clipped surrogate objective. TRPO's conjugate gradient computation also makes it incompatible with architectures that use dropout or parameter sharing between the policy and value networks, practical constraints that PPO eliminates. The clipped objective is:
\begin{equation}
L^{\text{CLIP}}(\theta) = \mathbb{E}_t \left[ \min\!\left( r_t(\theta) \hat{A}_t, \; \text{clip}(r_t(\theta), 1-\varepsilon, 1+\varepsilon) \hat{A}_t \right) \right]
\end{equation}
where $r_t(\theta) = \pi_\theta(a_t|s_t) / \pi_{\theta_{\text{old}}}(a_t|s_t)$ is the probability ratio and $\hat{A}_t$ is an estimator of the advantage function. The clipping removes the incentive for $r_t(\theta)$ to move outside the interval $[1-\varepsilon, 1+\varepsilon]$, effectively penalizing large policy updates without explicitly computing or constraining a divergence measure. PPO can be optimized by running multiple epochs of minibatch stochastic gradient descent on the same batch of trajectory data,\footnote{PPO typically runs 3--10 epochs of minibatch SGD per batch of trajectory data. The clipping threshold $\varepsilon = 0.2$ is standard. This multi-epoch reuse is unusual for on-policy methods (which traditionally discard data after one update) and is made safe by the clipping: once $r_t(\theta)$ exits $[1-\varepsilon, 1+\varepsilon]$, the gradient contribution from that sample is zeroed, preventing further policy drift on stale data.} making it substantially simpler to implement than TRPO.

In empirical comparisons, PPO outperformed A2C, TRPO, and the cross-entropy method on continuous control benchmarks and achieved the highest average reward on 30 of 49 Atari games among the methods tested. PPO became the most widely used policy optimization algorithm for large-scale RL applications, including the preference learning systems discussed in Section~\ref{section:rlhf}.

\subsection{AlphaGo and AlphaGo Zero (2016, 2017)}

The game of Go, with approximately $10^{170}$ legal positions and a branching factor of roughly 250, had remained beyond the reach of both brute-force search and hand-crafted evaluation functions. \citet{Silver2016} combined learning and planning by integrating multiple components: a policy network trained by supervised learning on 30 million positions from human expert games, a second policy network improved by self-play reinforcement learning, a value network trained to predict game outcomes from positions, and Monte Carlo tree search (MCTS) guided by these networks. AlphaGo defeated Lee Sedol, one of the strongest Go players in history, four games to one in March 2016.

\citet{Silver2017} removed human data entirely. AlphaGo Zero uses a single neural network $f_\theta(s) = (\mathbf{p}, v)$ that takes a board position $s$ as input and outputs a policy vector $\mathbf{p}$ over legal moves and a scalar value $v$ estimating the probability of winning. The training loop is self-contained: the program plays games against itself using MCTS guided by $f_\theta$, collects the MCTS-improved move probabilities $\boldsymbol{\pi}$ and the game outcome $z \in \{-1, +1\}$, then updates $f_\theta$ to minimize
\begin{equation}
\ell(\theta) = (z - v)^2 - \boldsymbol{\pi}^\top \log \mathbf{p} + c \|\theta\|^2
\end{equation}
where the first term is a value prediction loss, the second is a policy cross-entropy loss, and the third is $L_2$ regularization. After 72 hours of self-play on 4 TPUs,\footnote{AlphaGo Zero used a ResNet with 40 residual blocks (approximately 13 million parameters) and trained on 4.9 million self-play games. Each MCTS simulation performed 1{,}600 rollouts per move. The version that defeated Lee Sedol used a distributed system with 1{,}202 CPUs and 176 GPUs. The transition from AlphaGo to AlphaGo Zero demonstrated that removing human data and simplifying the architecture (single network instead of separate policy and value networks) improved both performance and training stability.} AlphaGo Zero surpassed all previous versions of AlphaGo, including the version that defeated Lee Sedol.

\citet{Igami2020} interprets the progression of game-playing programs through the lens of structural econometrics. Deep Blue is a calibrated value function: its 8,150 hand-tuned parameters encode expert knowledge directly, analogous to a fully calibrated structural model. Bonanza, a championship shogi program, estimates 50 million parameters from professional game data by minimizing prediction error on expert moves, a procedure Igami maps to Rust's nested fixed-point (NFXP) algorithm for maximum likelihood estimation of dynamic discrete choice models. AlphaGo combines conditional choice probability (CCP) estimation with forward simulation, connecting it to the Hotz-Miller two-step tradition. The progression from calibration to full-solution estimation to two-step estimation mirrors the methodological evolution of structural econometrics itself, a parallel that strengthens the bridge between reinforcement learning and economics that this survey is built around.

\subsection{Toward a Unified Framework}

The common object underlying all of these developments is the Bellman operator and its contractive properties. Thorndike's Law of Effect and the Rescorla-Wagner prediction error are the behavioral precursors of the TD error $\delta_t = r_{t+1} + \gamma V(s_{t+1}) - V(s_t)$. Samuel's weight update is TD(0) with linear function approximation. Q-learning is stochastic value iteration; natural policy gradient is soft policy iteration. DQN and AlphaGo stabilize function approximation through experience replay, target networks, and tree search. TRPO and PPO constrain the policy update to remain within a trust region. The difference between planning and learning is whether the Bellman operator is applied with known transitions or sampled ones. Chapter~\ref{section:planning_learning} develops this unified framework in full.
